\section{Threats to validity}
\label{sec:threats}

\paragraph{Construct validity.} Latency is attributed to layers by timestamps taken inside the Java processes.
Work that happens before the servlet filter runs (TCP accept, HTTP parsing in Tomcat's poller, waiting for a
worker thread) is not visible to the filter. On an inner hop it lands in the caller's \emph{net} layer, and on
the edge hop in the \emph{edge} layer. Under saturation both layers therefore contain queueing as well as
transport, and we interpret them that way. Application time is a residual of the handler's duration, so it also
absorbs JVM effects such as garbage-collection pauses that fall inside a handler. Throughput in the closed-loop
sweeps counts only successful responses. In~D part of the kernel's network processing for a request runs in
the sidecar's cgroup rather than the application's, which is consistent with D's edge application container
needing less CPU per request than C's. D's application budget therefore buys slightly more application work than
C's. On the Apple-silicon host, vCPUs run on both performance and efficiency cores, so the same CFS quota buys
different amounts of work depending on where a vCPU is scheduled. Throttling measured there may overstate what
the same quota causes on uniform server cores. The USL is a phenomenological model; a good fit does not identify
the mechanism behind $\alpha$ and $\beta$, which is why Section~\ref{sec:results} pairs every fit with the
measured CPU utilisation and throttling of the containers involved.

\paragraph{Internal validity.} The configurations run the same binary, receive the same request sequence and are
checked for identical responses, so implementation differences are excluded by construction. Every cell is a
fresh deployment and the configuration order is shuffled per repetition, which guards against drift of the
host. Three settings of the stack were changed from their defaults after calibration: a 1\,GiB memory floor per
application pod, a cap of 1{,}000 connections per Tomcat process, and server-side keep-alive without a request
limit. Each applies to all configurations and is documented with its reason in Section~\ref{sec:budget}. The budget is enforced as hard per-pod limits. Without limits, a decomposed configuration's busy edge service
could borrow the CPU its idle neighbours leave unused, so the fragmentation penalty under the equal split would
shrink, while the per-request CPU cost of distribution would remain. With the CPU per request of
Table~\ref{tab:cost} (\xo{cpumsref/browse/C/total} against \xo{cpumsref/browse/A/total}\,ms for browse), even a
pooled six-core budget would serve C less than half as many requests as the monolith. The equal CPU split is a deliberate choice, not
the only defensible one: an operator
who sizes each service to its demand would give the edge service of~C more than half a core. The results on
fragmentation (Section~\ref{sec:results}) are therefore results about equal static allocation. The two further
treatments of Section~\ref{sec:treatments} test the alternatives. Each has its own bias. The proportional shares are
fitted to the very workloads they are evaluated on, which favours the decomposed configurations. The autoscaling
results depend on the controller's settings and on the step length. With 180\,s levels, a new replica's JVM
warm-up falls inside the measurement, which is part of what elasticity costs. A longer level, a lower utilisation
target or warmed replicas would shift the sustained rates.

The procedure has a fixed order. Rate and concurrency levels always ascend, browse always precedes checkout on the
same deployment, and the closed-loop sweep follows the open-loop sweep. Level is therefore confounded with time
since deployment, and a decomposed configuration starts its closed-loop sweep after its overloaded open-loop
levels. The recovery probe limits this carry-over but cannot rule it out, and it may contribute to the dip in the
closed-loop curves of B and C at moderate concurrency. Database state also changes during a cell. Orders and
reservations are append-only and the ten popular stock rows are updated in place, so the databases grow and
accumulate dead tuples, most in the monolith's single database, which serves every write; autovacuum may run
during measurement. All PostgreSQL instances keep their data on volumes of their node with \texttt{fsync} enabled,
and in~C a checkout commits in three databases that can share one disk. In the proportional campaign, B, C and~D
are redeployed for checkout with the checkout shares, whereas A, whose allocation is the same for both workloads,
runs checkout on the deployment that served browse. Finally, some data were excluded. In the proportional campaign
three open-loop steps of one cell (C, browse at 600, 800 and 1{,}000 requests per second) were dropped because the
host slept. An earlier local campaign that a host suspension interrupted was discarded entirely, as was the first
cloud campaign (Section~\ref{sec:procedure}). Their logs are archived with the data.

\paragraph{The measurement environment.} All nodes of the kind cluster, including the load generator's, share
one host's CPUs. The load generator was not CPU-limited, and the system under test could not take more than its
CPU limits, but the two still share caches, memory bandwidth and the hypervisor's scheduler. Inter-node traffic
crosses a Linux bridge inside one virtual machine rather than a physical network, so the \emph{net} layer here is
a lower bound on what the same topology costs across real hosts. We also observed a low-load penalty that is a
property of the host rather than of any configuration: in the monolith, browse p99 was \xo{p99/browse/A/50}\,ms at 50 requests per second but
\xo{p99/browse/A/800}\,ms at 800. A separate control run ruled out JIT warm-up as the cause. The penalty persisted
after a closed-loop warm-up of about 580{,}000 requests (p99 4.0\,ms at 50 requests per second against 0.85\,ms at
800), and it reappeared on returning to 50 requests per second after the high-rate step (4.8\,ms). We attribute it to CPU idle-state exits and thread wake-ups when requests arrive 20\,ms apart.
It affects every configuration, but configurations with more hops wake more threads per request, so it inflates
their low-load ratios in Table~\ref{tab:h1}. At the lowest rate the H1 ratios were
\xo{h1/min}--\xo{h1/max}, but at 200 browse requests per second they were \xo{ratio/browse/B/200/est} (B),
\xo{ratio/browse/C/200/est} (C) and \xo{ratio/browse/D/200/est} (D). The size of the low-load penalty therefore
depends on the reference rate, while its sign does not. On bare metal it can be controlled by pinning the
CPU frequency governor and limiting idle states, which is part of the protocol for the remaining environments.

Low-load tail latency also drifted between campaigns. The monolith, deployed identically, had a browse p99 at
50 requests per second of \xo{p99/browse/A/50}\,ms in the first local campaign and \xo{prop/p99/browse/A/50}\,ms in
the second. The H1 ratios of one campaign are therefore more precise than their reproducibility across campaigns
on this host.

The cloud environment removes the shared host but adds variation the local one does not have. The instances are
on shared tenancy, so other customers' workloads can affect caches and the network. The VPC network's latency
varies over time. Instances of one type can also differ slightly in performance. Every node sat in one
availability zone of one region, so the results do not include cross-zone traffic. The default scheduler's
placement also decides which calls cross nodes. It is recorded per cell but not varied, so placement is part of
each configuration in this environment rather than a controlled factor.

\paragraph{External validity.} The application is small and synthetic, written in one language and one
framework. Its per-request business logic is cheap, so the fixed costs of distribution (serialisation, HTTP,
network) weigh more than they would in a service that spends milliseconds computing. The service-demand treatment
adds 1\,ms of CPU to every module operation to show how the conclusions shift as that ratio changes. Each run also
exercises one operation, so the modules it does not use are idle by construction. A mixed workload would spread
demand across services and reduce the fragmentation penalty of the equal split. The proportional allocation, tuned
to each operation, cannot be applied to such a mix as it stands. Communication is
JSON over HTTP/1.1; binary protocols such as gRPC over HTTP/2 reduce serialisation and connection costs and
could change the magnitudes, although not the direction, of the per-hop results. All databases are PostgreSQL
instances co-located in the same cluster. The mesh is Istio~1.31 with sidecars; sidecar-less data planes such as
Istio Ambient are reported to cost far less~\cite{BremlerBarr2025} and were not measured. Both environments use
arm64 processors, one instance family and one cloud provider. x86 servers, bare metal and other providers,
which the same scripts also target, remain untested.

\paragraph{Conclusion validity.} With \xo{reps} repetitions per configuration, intervals are wide for effects
that are small relative to the variation between deployments. We report intervals rather than binary
significance wherever possible, correct the H1 family for multiple comparisons, and treat the repetition, not
the request, as the unit of analysis. The raw results of every cell are published so that other analyses can be
run on them.
